\documentclass[a4paper]{article}
\usepackage[pdftex]{graphicx}
\usepackage[utf8]{inputenc}
\usepackage{enumerate}
\usepackage{siunitx}
\sisetup{locale=DE}
\usepackage{href-ul}
\hypersetup{
	colorlinks=true,
	linkcolor=blue,
	urlcolor=blue}
\usepackage{icomma}
\usepackage{amssymb}
\usepackage{geometry}
\geometry{a4paper, top=15mm, left=15mm, right=15mm, bottom=15mm,
	headsep=10mm, footskip=12mm}

\begin{document}
	
	\begin{enumerate}[1.]
		
		{\bf \item What is meant by $\dots$}
		
		\begin{enumerate}[a)]
			\item $\dots$ damped oscillation
			\item $\dots$ Amplitude of an oscillation
		\end{enumerate}
		
		%Exercise 2
		{\bf \item A pendulum makes 20 oscillations in 5 seconds. Find the frequency of the pendulum.}
		
		%Task 3
		{\bf \item A spring pendulum (spring hardness $D=\SI{1.9}{\newton\over\meter}$, mass $m=\SI{200}{\gram}$) is located at time $t =0$ at the lower reversal point.}
		
		\begin{enumerate}[a)]
			\item Calculate the period of oscillation of this spring pendulum to the nearest tenth of a second. Also specify the frequency.
			\item Use a sketch to explain where the swinging pendulum body is
			$t_1=\SI{3,5}{\second}$ is located and in which direction it is moving.
		\end{enumerate}
		
		%task 4
		{\bf \item For a spring pendulum (pendulum mass $m$, spring hardness $D$), the oscillation period should be halved.} Name two different possibilities.
		
		
		%Task 5
		{\bf \item The oscillation of a spring pendulum (spring hardness $\SI{0.49}{\newton\over\meter}$, amplitude $\SI{12}{\centi\meter}$)}
		starts at the lower reversal point at time $t=0$ and reaches the upper reversal point two seconds later.
		
		\begin{enumerate}[a)]
			\item What is the oscillation period?
			\item How big is the attached mass?
			\item How large is the restoring force and the acceleration at the upper reversal point?
			\item At which first three points in time is the speed of the pendulum body greatest?
		\end{enumerate}
		
		{\bf \item For another harmonic oscillation} with the amplitude $ \hat{x}=\SI{10}{\centi\meter} $ you measure an oscillation period of $\SI{0.4}{\second} $ Explain what frequency this oscillation has at an amplitude of $ \hat{x}=\SI{20}{\centi\meter} $.
		%Task 3
		{\bf \item From the diagrams below, determine the amplitude, wavelength, period of oscillation} and the frequency and speed of propagation of the mechanical wave.
		
		\includegraphics[width=12 cm]{sinus95.pdf}
		
	\end{enumerate}
	
	\fbox{
		\begin{minipage}{0.5\textwidth}
			For the solution, please \href{https://www.okuyakl.de/phys/p10pebL408/le408.pdf}{click here} or scan the QR code.\\
			Additional worksheets are available at
			
			\href{http://www.okuyakl.com}{www.okuyakl.com}
		\end{minipage}
		\hfill
		\begin{minipage}{0.4\textwidth}
			\includegraphics[width=1.5 cm]{../../viecher/zwe03}
			\includegraphics[width=3 cm]{qre408}
			\includegraphics[width=2 cm]{../../viecher/afanticon1}
			
	\end{minipage}}
	
\end{document}